\clearpage
\section{Project development}

This section describes the decisions taken during the development of the project. The previous chapter defined the methodology. In this chapter, the focus is how the requirements were translated into the structure of the tool: why the exploration was divided into two levels, how the different parts are connected, and the scope that guided the implementation.

The contribution of the thesis is Talos, a Python tool for hierarchical design space exploration of DNN accelerators. Talos is intended as an architectural exploration tool: its purpose is to compare accelerator alternatives for a workload and a set of design constraints, first at an abstract architectural level and later using characterized hardware IP blocks. Generating synthesizable RTL or producing a single final accelerator is outside the intended scope.

\subsection{Tool Overview}
Talos is programmed in Python. The idea behind the software design is to make Talos easy to work on and modifiable without having to have broad knowledge about the full implementation. This is important because the tool is intended to be extended with new evaluation models, optimization objectives, workloads, and IP libraries.

The intended state of Talos is that a user can provide a DNN workload, a set of optimization objectives and constraints, and an IP library, and obtain a traceable set of feasible architectural alternatives. The tool is therefore designed to support architectural decisions and comparisons, rather than to hide the exploration process behind a single generated design.

Talos is composed of two different stages, or levels. The first level is the Architectural Exploration. This level reduces the design space by evaluating candidates at an architectural level, without stopping to consider the possible IPs that could be used.

The second level is the IP-Based Hardware Evaluation. This step evaluates the candidates selected by the previous level taking into account the different IPs available.

The intended output is a set of candidate architectures with specific IP choices and comparable metrics. Each candidate is represented internally as an implemented accelerator composed of selected hardware IP blocks. This representation does not correspond to synthesizable RTL; it is the final abstraction needed by Talos to evaluate and compare architectural alternatives.

The adopted flow can be summarized as follows:

\begin{enumerate}
\item A genome describes an accelerator candidate at an architectural level.
\item The genome is decoded into an internal architecture configuration.
\item This configuration is translated into a ZigZag-compatible accelerator description.
\item ZigZag is used to evaluate the candidate in terms of energy, latency, and area-related cost.
\item The architectural candidate can be converted into an abstract accelerator representation.
\item The second level selects compatible hardware IP blocks for the abstract components.
\item The implemented accelerator is evaluated with a workload-aware PPA model and checked against the user constraints.
\end{enumerate}

This structure reflects the main idea of the project: keeping architectural exploration separated from hardware IP instantiation while maintaining a path between both stages.

\subsection{Abstract Architecture Representation}
In order to evaluate the different architectures, an abstract representation that does not contain specific IPs is needed. Since ZigZag is used in the first level, the decision was to keep this representation close to ZigZag's hardware model.

This thesis will not go in depth about all the technical details of their work, but a brief overview is useful to understand how Talos represents the accelerator before moving to the IP-based evaluation stage.

In Talos, this abstract representation appears at two different points. First, in Level 1, an architecture is represented through a compact genome and then decoded into a configuration that can be exported to ZigZag. Second, when moving to Level 2, the same candidate is translated into an internal \texttt{AbstractAccelerator}. This object contains abstract components such as the PE array, register files, and global buffer, but it does not yet select concrete IP implementations.

This separation is important because the first level only needs to know what the architecture looks like, while the second level needs to know which hardware blocks could implement it.

\subsection{Hardware architecture representation in ZigZag}

ZigZag models the hardware architecture of a DNN accelerator using a hierarchical representation that starts from the smallest compute primitive and progressively builds the full accelerator. At the lowest level, the \textit{operational unit} represents the basic compute element, typically a multiplier, and is characterized by attributes such as energy and area. Multiple operational units are then organized into an \textit{operational array}, whose dimensions define the available spatial parallelism.

The memory system is described through \textit{memory instances}, each one parameterized by capacity, read and write bandwidth, access cost, area, and number of ports. These memory instances are connected into a \textit{memory hierarchy}, which captures both storage properties and the connectivity between memory levels and the operational array. A relevant concept in this representation is that of \textit{memory operands}, which decouples the hardware memory organization from the algorithmic tensors and later links them through the mapping specification.

ZigZag also models how memories are spatially distributed with respect to the compute array by means of the \texttt{served\_dimensions} attribute. This allows the framework to represent different organizations, from private memories attached to each operational unit to shared buffers serving one or more array dimensions. In addition, the hierarchy explicitly models the allowed data movements between adjacent levels, including transfers from higher to lower levels and vice versa.

Finally, the combination of the operational array and the memory hierarchy forms a \textit{core}, and this core is wrapped into the top-level \textit{accelerator} object. This representation is especially relevant because it provides an explicit and structured description of the compute resources, memory organization, and communication patterns of the accelerator, which are key aspects during design space exploration and mapping evaluation.

The generated ZigZag accelerator follows a simple and fixed memory structure. It contains a PE array, one register file for input activations, one register file for weights, one register file for outputs and partial sums, a shared global buffer, and an external DRAM level.

The operand links used by Talos are:

\begin{itemize}
\item input activations are mapped to \texttt{I1};
\item weights are mapped to \texttt{I2};
\item outputs are mapped to \texttt{O}.
\end{itemize}

The external DRAM is kept in the ZigZag description because it is needed to model off-chip traffic. However, it is not treated as part of the accelerator area. In the current model, the DRAM has zero area and a fixed bandwidth. This decision avoids letting the optimizer increase off-chip bandwidth without modelling the physical cost of the external memory interface.

\subsection{Level 1: Architecture Exploration}
The first level of Talos performs architectural exploration. Its objective is to generate and evaluate accelerator candidates without selecting concrete hardware IP blocks yet.

The architecture is encoded as a discrete genome. Each gene is an index into a catalogue of possible values. This representation was selected because it is simple, compatible with evolutionary optimization, and avoids generating arbitrary values that do not correspond to the current design space.

The current Level 1 genome contains seven genes, shown in Table~\ref{tab:level1_genome}.
\todo[inline]{Revisitar genoma, ya que este register file puede estar dentro del PE}
\begin{table}[h]
\centering
\begin{tabular}{lll}
\hline
Gene & Meaning & Role \\
\hline
\texttt{pe\_x\_code} & PE array size in the first dimension & Explored \\
\texttt{pe\_y\_code} & PE array size in the second dimension & Explored \\
\texttt{rf\_size\_code} & Register file capacity & Explored \\
\texttt{rf\_bw\_code} & Register file bandwidth & Explored \\
\texttt{gb\_size\_code} & Global buffer capacity & Explored \\
\texttt{gb\_bw\_code} & Global buffer bandwidth & Explored \\
\texttt{gb\_served\_dims\_code} & Dimensions served by the global buffer & Explored \\
\hline
\end{tabular}
\caption{Current Level 1 genome used by Talos.}
\label{tab:level1_genome}
\end{table}

The DRAM bandwidth was intentionally removed from the genome. In previous iterations it was considered as an architectural gene, but this introduced an unrealistic optimization possibility: the optimizer could improve latency by selecting a higher off-chip bandwidth without paying any area or power cost for the corresponding interface. In the current version, DRAM bandwidth is therefore treated as a fixed platform parameter.

After decoding the genome, Talos builds an architecture configuration. This configuration is then used to generate the accelerator description consumed by ZigZag. The Level 1 evaluation produces metrics such as energy, latency, and area-related information. When an area value is not directly available, Talos uses a simple fallback area model to keep the flow executable.

The first level also includes an objective adapter. This component allows the same raw evaluation result to be exposed as different optimization objectives, such as latency, energy, area, energy-delay product, energy-area product, or area-latency product. This avoids hard-coding a single objective into the evaluator and makes the optimization process easier to extend.

\subsection{IP representation and hardware pool}
The second level of Talos requires a representation of the available hardware blocks. For that reason, an IP pool was implemented.

Each IP block represents a possible hardware component that can be selected to implement part of the accelerator. The current IP representation contains the following information:

\begin{itemize}
\item IP name;
\item IP type;
\item area;
\item idle and active power characterization;
\item delay;
\item throughput;
\item reference frequency;
\item optional capacity;
\item optional bandwidth;ion voltage, temperatu
%\item characterizatre, and process corner.
\end{itemize}

The IP type is used to match hardware blocks with abstract accelerator components. For example, a register file component must be matched with a compatible memory IP, and a PE array component must be matched with a compatible compute IP.

Capacity and bandwidth are optional fields because not all IP types require them in the same way. However, when they are present, they are used during compatibility checking. For example, if an abstract memory component requires a minimum capacity or bandwidth, an IP block that does not satisfy those requirements is rejected.

The IP pool is loaded from a structured description file. This makes it possible to modify the set of available IPs without changing the code. The intended use is to populate this pool with values obtained from a coherent characterization flow, such as synthesis, memory compilers, CACTI, or other estimation tools.

During the early stages of development, synthetic characterization values were used only to implement and validate the exploration flow. Once the flow was stable, the IP pool used for the final evaluation was populated with real characterization values.

\subsection{Level 2: IP-Based Hardware Evaluation}
The second level evaluates an architectural candidate using the available hardware IP pool. The input to this level is not the original genome, but an abstract accelerator generated from the Level 1 architecture.

This abstract accelerator is composed of several abstract components. In the current implementation, these components usually include:

\begin{itemize}
\item the PE array;
\item one register file for input activations;
\item one register file for weights;
\item one register file for outputs and partial sums;
\item one shared global buffer.
\end{itemize}

The external DRAM is intentionally not instantiated as a Level 2 IP block. It is part of the memory environment used during ZigZag evaluation, but it is not considered an on-chip accelerator component. This is consistent with the area model used in Level 1.

The Level 2 genome is generated dynamically from the abstract accelerator and the IP pool. Instead of having a fixed number of genes, each gene corresponds to one abstract component. The value of each gene selects one compatible IP implementation for that component.

This design allows the second level to adapt to different accelerator candidates. If the abstract accelerator changes, the Level 2 genome changes with it. This is useful because the number and type of components may evolve as the architectural representation becomes richer.

Once a Level 2 genome is decoded, Talos builds an implemented accelerator. This implemented accelerator contains the selected IP block for each abstract component. The evaluator then computes the physical metrics and checks whether the candidate satisfies the user constraints.

Area is accumulated across the selected components, taking into account the number of instances. Delay is approximated using the maximum delay among the selected blocks, throughput using the minimum throughput, and the implementation frequency using the lowest maximum frequency of its components.

For workload-aware energy and power, Talos reuses the mapping selected by ZigZag. For each layer, it considers its execution time, the number of active PEs, and the accesses to the register files and global buffer. These values are combined with the idle and active power characterization of each selected IP. Dynamic DRAM access energy is also included, although DRAM remains outside the on-chip area and IP selection. All selected IPs use a common characterization frequency and compatible voltage, temperature, and process corner so that the comparison remains coherent.

\subsection{Optimization Objective Interface}
The original objective of this part of the tool was to provide a modular fitness framework. In the current version, this idea is implemented as an objective interface that separates evaluation from objective selection.

This separation is useful because the evaluator may return several metrics, but the optimizer may not always use the same ones. For example, one experiment may optimize latency and energy, while another may optimize area and energy-delay product. Instead of modifying the evaluator for each case, Talos exposes different objective functions on top of the same evaluation result.

In Level 1, the available objectives include:

\begin{itemize}
\item latency;
\item energy;
\item area;
\item energy-delay product;
\item energy-area product;
\item area-latency product.
\end{itemize}

In Level 2, the objectives are based on the values obtained after selecting IP blocks. They include area, energy per inference, average workload power, delay, and inverse throughput. Inverse throughput is used because the optimizer follows a minimization convention.

A cache mechanism is also included so that the same candidate does not need to be evaluated multiple times when several objectives are requested from the same result. This is especially useful when the evaluation step is expensive.

\subsection{Design Flow}
The current Talos flow starts with the Level 1 architectural genome. This genome is decoded into an architecture configuration and exported to ZigZag. ZigZag evaluates the candidate architecture and returns the architectural metrics used by the optimizer.

After this, the candidate can be imported into Level 2. The importer converts the architecture into an abstract accelerator representation. This representation is then compared against the IP pool to find compatible implementations for each component.

The Level 2 genome is created from those compatibility lists. Each candidate genome selects one IP block per abstract component. The selected IPs are used to build an implemented accelerator, and the Level 2 evaluator computes its physical and workload-aware values.

The complete flow can therefore be seen as a hierarchical exploration:

\begin{enumerate}
\item Level 1 explores architectural candidates.
\item ZigZag evaluates the mapping and memory behaviour of those candidates.
\item Selected candidates are converted into abstract accelerator components.
\item Level 2 explores possible IP implementations.
\item A workload-aware PPA model evaluates the implemented candidate.
\item User constraints discard architectures that are not feasible.
\item The feasible alternatives and their metrics are exported for comparison.
\end{enumerate}

Level 1 uses NSGA-II through \texttt{pymoo} and a fixed-length architectural genome. Level 2 uses a dynamic IP-selection genome and can use either NSGA-II or exhaustive enumeration. Exhaustive exploration is useful when the compatible IP pool is small, while NSGA-II preserves the same flow when the number of combinations becomes too large. The output is a set of feasible alternatives, including non-dominated candidates, rather than a single architecture selected by the tool.

\subsection{Challenges and Problems Encountered}
One of the main challenges was keeping the meaning of the metrics 
coherent between both levels. Level 1 evaluates an abstract 
architecture with ZigZag, while Level 2 evaluates concrete IP choices.
When the available IP combinations are introduced, an architecture
that appears less promising at the abstract level may produce a better
physical implementation, while a well-ranked abstract architecture
may offer less favourable IP combinations. For this reason, the
connection between levels was designed to preserve several candidate
architectures instead of passing only one apparent optimum.

Another problem was the distinction between energy and power. 
In an early version of the flow, ZigZag energy was used as a proxy 
for a power-oriented objective. This produced misleading 
comparisons because a large architecture may consume less 
energy by finishing earlier while still requiring more average 
power. The decision was to keep energy and power as different 
quantities and to make the Level 2 model workload-aware, 
using the execution time and activity of each layer.

Frequency also required special attention. The maximum frequency 
of an IP indicates whether it can operate at a given point,
but it is not automatically the frequency at which its power 
was characterized. Combining power values from one frequency 
with execution time calculated at another produced inconsistent 
energy estimations. Talos therefore evaluates workload time at 
the common reference frequency of the selected IP characterizations
and uses maximum frequency as a feasibility condition.

The treatment of DRAM exposed a similar modelling issue. 
Allowing the optimizer to change external bandwidth without 
assigning a physical cost created an unrealistic path to lower latency. Fixing DRAM bandwidth as a platform parameter avoided this problem. Later, the dynamic energy associated with off-chip accesses also had to be propagated from the ZigZag mapping, because ignoring it favoured small on-chip architectures with high external traffic.

The optimization algorithms introduced practical problems 
as well. Although the Level 1 genome is discrete, generic 
genetic operators can produce floating-point individuals 
that decode into repeated architectures and reduce population 
diversity. The genome operators were therefore kept integer 
throughout the search. In Level 2, NSGA-II could miss feasible 
combinations when constraints were strict, even when the small 
IP pool allowed all combinations to be checked. An exhaustive strategy was added for this case, leaving NSGA-II for larger pools where enumeration is no longer practical.

Finally, the accuracy of the exploration depends on the quality of 
the external models. ZigZag does not always provide the required 
area information, automatic memory characterization is 
not equally portable in every environment, and a complete real 
IP library was not available during the early stages of development. 
Talos keeps the source of the estimates explicit: synthetic values 
were used while developing the flow, while the final evaluation 
used the pool obtained from real characterization.

\subsection{Validation of the Design Decisions}
The validation focuses on functional correctness and integration between modules. The goal is not to claim final physical accuracy, but to verify that the adopted structure works as an architectural exploration tool.

The implemented tests cover several parts of the flow:

\begin{itemize}
\item validation of IP block definitions;
\item loading of IP pools from YAML files;
\item compatibility checks between abstract components and IP blocks;
\item rejection of invalid candidates when no compatible IP is available;
\item generation and decoding of Level 2 genomes;
\item construction of implemented accelerators;
\item computation of area, delay, throughput, workload energy, and average power;
\item application of area, power, latency, and frequency constraints;
\item consistency between characterization frequency, workload time, energy, and power;
\item conversion from Level 1 architectures to Level 2 abstract accelerators.
\end{itemize}

In addition to unit tests, smoke tests are used to verify that the modules can be connected in practice. One smoke test creates an architecture from the Level 1 representation, converts it into an abstract accelerator, loads an example IP pool, generates the corresponding Level 2 genome, and evaluates the resulting implementation.

These tests check that Talos can represent, transform, and evaluate accelerator candidates through both levels. They also check an important modelling decision: the DRAM is preserved in the ZigZag architecture for off-chip traffic modelling, but it is not imported as an IP block in Level 2.

\subsection{Current Limitations}
The current version of Talos should be understood as a working prototype of an architectural exploration tool. Some parts are intentionally simplified and define how its evaluations should be interpreted.

First, the Level 1 area model is still preliminary. ZigZag provides useful architectural evaluation, but area extraction may not always be directly available or sufficiently detailed. For that reason, Talos includes a fallback area model. This is useful for exploration, but it is not a replacement for a technology-accurate area model.

Second, the Level 1 design space is defined through a fixed catalogue of discrete values. This keeps candidate generation valid and manageable, but it also makes the exploration rigid and may leave relevant architectural values outside the search.

Third, the Level 2 evaluator does not yet model pipeline effects, interconnect and clock-tree cost, memory arbitration, control logic, DRAM idle power, or the physical memory interface. These aspects are necessary for a more accurate hardware evaluation. Although the final IP pool uses real characterization values, its accuracy is limited by the components and operating conditions covered by that characterization.

Finally, the separation between levels makes the exploration manageable by applying the more detailed IP evaluation only to a representative set of architectural candidates. Preserving a diverse set of Level 1 candidates allows Level 2 to compare different architectural alternatives and concrete IP combinations without requiring an exhaustive joint exploration of both spaces.

\subsection{Summary of Project Decisions}
The development of Talos was guided by the decision to build a modular architectural exploration tool rather than a generator of one final accelerator.

The main decisions taken during the project were:

\begin{itemize}
\item to divide the exploration into an architectural level and an IP-based level;
\item to reuse ZigZag for workload mapping and architectural evaluation;
\item to encode Level 1 as a discrete genome and keep DRAM as a fixed external platform resource;
\item to connect both levels through an abstract accelerator representation;
\item to describe physical alternatives through a characterized IP pool and a dynamic Level 2 genome;
\item to keep evaluation metrics separate from optimization objectives;
\item to evaluate energy and power from workload activity and coherent IP characterization;
\item to use NSGA-II for large spaces and exhaustive Level 2 exploration when the IP pool is small;
\item to preserve several feasible alternatives and report their trade-offs and constraint status.
\end{itemize}

The intended state of Talos is therefore a tool that helps a designer explore and compare DNN accelerator architectures with traceable assumptions. 
Its endpoint is a set of feasible architectural alternatives with concrete IP choices and comparable metrics; RTL generation and sign-off physical 
estimation remain outside this scope.

%\include{codi_tempdistrib}
